\chapter{Proposed Method}
\label{ch:proposed_method}

This chapter presents our approach to solving the constrained 3D bin packing problem. Rather than describing the system by component type, we follow the development journey: starting with the geometric foundation of 3D nesting, building up a reinforcement learning framework to learn packing strategies, designing a neural architecture to process spatial information, and finally employing evolutionary methods to optimize the architecture itself.

\section{Introduction}

From the problem analysis in Chapter~\ref{ch:problem_analysis} and the survey of existing methods in Chapter~\ref{ch:state_of_art}, several key insights emerged:

\begin{itemize}
    \item The problem is NP-hard with no polynomial-time optimal solution~\cite{garey1979computers}.
    \item Supervised learning approaches require labeled datasets of optimal solutions, which are intractable to compute at scale.
    \item Heuristic methods struggle with the combinatorial complexity of multiple constraints (incompatibility, affinity, relative positioning).
    \item Spatial reasoning is crucial---placement decisions must account for local geometry and long-term packing potential.
\end{itemize}

Our approach combines geometric heuristics with deep reinforcement learning. We first establish a robust geometric foundation using Empty Maximal Spaces (EMS) for collision-free candidate generation. We then layer reinforcement learning on top to learn optimal selection among candidates. Finally, we employ neural architecture search to discover effective network designs automatically.

\section{Building the Foundation: 3D Nesting}

The first challenge is the geometry itself: how do we efficiently generate valid placement positions in a 3D container?

\subsection{Empty Maximal Spaces}

As discussed in Section~\ref{sec:geometric_heuristics}, two main geometric heuristics exist: Extreme Points (EP)~\cite{crainic2008extreme} and Empty Maximal Spaces (EMS)~\cite{crainic2012packing}. We choose EMS for several critical advantages:

\begin{enumerate}
    \item \textbf{Collision-Free by Construction}: Each EMS represents a maximal empty rectangular volume. Any item placement within an EMS is guaranteed not to collide with existing items, eliminating the need for collision checking at decision time.

    \item \textbf{Efficient Updates}: When an item is placed, we only need to update the EMS set once---splitting intersected spaces and removing dominated ones. With Extreme Points, collision checks must be performed for every candidate action at every decision step.

    \item \textbf{Bounded Action Space}: EMS naturally prunes infeasible placements. If an item is too large to fit in any EMS, it is automatically excluded from the action set.

    \item \textbf{Quality Metric}: The size and shape of available EMS spaces provides a geometric quality metric that can guide learning (detailed in Section~\ref{sec:reward_shaping}).
\end{enumerate}

The computational efficiency is particularly important for reinforcement learning, which requires millions of environment interactions during training.

\subsection{Action Space}

With EMS providing the geometric foundation, we define the action space. Each action represents a complete placement decision:

\begin{equation}
a = (\text{bin\_idx}, \text{item\_idx}, \text{ems\_idx}, \text{rot\_idx})
\end{equation}

At each decision step, the environment enumerates all valid combinations: for each bin, for each remaining item, for each EMS that can fit the item, for each of the 6 possible rotations. This generates anywhere from dozens to hundreds of candidate actions depending on problem state.

To enable learning, each action must be represented numerically. We encode actions as feature vectors capturing geometric properties (positions, dimensions, fit metrics), stability indicators (support, center of mass), and contextual information (bin utilization, remaining capacity). These features provide the neural network with information about both immediate placement quality and global packing state.

\subsection{Heightmap Representation}

Scalar features alone cannot capture spatial structure. Consider two placements with identical utilization and fit metrics, but one places an item in a valley (creating a flat surface for future items) while the other creates a protruding tower (fragmenting the space). To address this, we need a spatial representation.

We represent each bin as a heightmap: a 2D grid where each cell stores the maximum occupied height at that $(x, y)$ location. For each candidate action, we extract a local patch centered on the placement position, providing a window into the surrounding geometry.

% \begin{figure}[h]
% \centering
% \includegraphics[width=0.8\textwidth]{figures/heightmap_example.png}
% \caption{Heightmap representation showing different placement scenarios: (left) corner placement, (middle) valley filling, (right) fragmentation.}
% \label{fig:heightmap}
% \end{figure}

This local spatial information reveals:
\begin{itemize}
    \item \textbf{Local Geometry}: Whether the location is in a corner, against a wall, in a valley, or on a plateau
    \item \textbf{Support Patterns}: How well the new item will be supported
    \item \textbf{Future Packing Potential}: Whether placement creates a flat surface or fragmentation
\end{itemize}

The patch must be large enough to capture surrounding context, yet small enough to remain computationally tractable during training.

\section{Learning to Pack: Reinforcement Learning}

With the geometric foundation established, the question becomes: given dozens or hundreds of valid placements generated by EMS, which one should we choose?

\subsection{Why Reinforcement Learning?}

Traditional supervised learning requires labeled training data---optimal solutions for problem instances. But computing optimal solutions to 3D bin packing is intractable due to NP-hardness~\cite{garey1979computers}. Even near-optimal heuristic solutions are expensive to generate and may not generalize.

Reinforcement learning learns directly from interaction~\cite{sutton2018reinforcement}. The agent learns through trial and error, receiving rewards based on packing quality. This eliminates the need for pre-computed solutions.

Furthermore, bin packing is inherently sequential: each placement affects future possibilities. This aligns naturally with the Markov Decision Process (MDP) framework underlying RL. The agent learns to maximize cumulative reward over the entire packing sequence, not just immediate placement quality.

\subsection{Deep Q-Networks}

We employ Deep Q-Networks (DQN)~\cite{mnih2015human}, a value-based reinforcement learning algorithm. DQN learns a Q-function $Q(s, a)$ estimating the expected cumulative reward of taking action $a$ in state $s$.

Key advantages for our problem:
\begin{itemize}
    \item \textbf{Sample Efficiency}: Experience replay allows learning from past experiences multiple times
    \item \textbf{Stability}: Target networks provide stable learning targets
    \item \textbf{Generalization}: Deep networks generalize across similar states
\end{itemize}

We adopt the Double DQN variant~\cite{van2016deep} to address Q-value overestimation. Standard DQN uses the max operator for both action selection and evaluation:

\begin{equation}
Q(s, a) \leftarrow r + \gamma \max_{a'} Q_{\text{target}}(s', a')
\end{equation}

Double DQN decouples these operations:

\begin{equation}
Q(s, a) \leftarrow r + \gamma Q_{\text{target}}(s', \argmax_{a'} Q_{\text{online}}(s', a'))
\end{equation}

This reduces overestimation, particularly important where many actions have similar immediate rewards but different long-term consequences.

\subsection{State Representation}

The state representation must capture both global progress and local geometry. We encode:

\paragraph{Global Features} (scalars):
\begin{itemize}
    \item Current bin utilization, bin index, total bins
    \item Number of items placed and remaining
    \item Bin dimensions and weight capacity
\end{itemize}

\paragraph{Item Features} (per-item vectors for remaining items):
\begin{itemize}
    \item Dimensions, weight
    \item Incompatibility flags, affinity group ID
    \item Relative positioning constraints
\end{itemize}

Item features are processed through a permutation-invariant encoder to handle variable numbers of remaining items.

\subsection{Reward Design}
\label{sec:reward_shaping}

Bin packing presents a sparse reward problem: the primary objective (minimizing bins) only provides feedback at episode end. Intermediate placements receive no signal about quality, creating severe learning challenges:

\begin{itemize}
    \item \textbf{Credit Assignment}: Which of hundreds of decisions led to needing an extra bin?
    \item \textbf{Slow Learning}: Random exploration rarely produces good solutions
\end{itemize}

To address this, we employ potential-based reward shaping~\cite{ng1999policy}, which augments the sparse environment reward with auxiliary signals:

\begin{equation}
r_{\text{shaped}}(s, a, s') = r_{\text{env}}(s, a, s') + \gamma \cdot \Phi(s') - \Phi(s)
\end{equation}

This formulation has a critical theoretical property: it guarantees the optimal policy under shaped reward is identical to the optimal policy under the original reward. We can guide learning without biasing the final solution.

Our potential function $\Phi(s)$ combines two intuitions:

\begin{itemize}
    \item \textbf{Bin Utilization}: Reward filling bins efficiently, measured by volume utilization
    \item \textbf{Space Quality}: Reward maintaining well-shaped empty spaces that can accommodate future items. EMS provides a natural metric---the largest remaining spaces indicate packing potential. By focusing on the largest spaces rather than total empty volume, we discourage fragmentation into many small, unusable gaps.
\end{itemize}

At episode termination, the agent receives reward based on the number of bins used, with strong penalties for exceeding the minimum. The shaped intermediate rewards guide the agent toward configurations that achieve this terminal objective.

\section{Neural Architecture}

Now that we have established what to learn (Q-function via DQN) and what signals guide learning (shaped rewards), we address how to represent the Q-function with neural networks.

\subsection{Processing Spatial Information: Heightmap CNN}

Heightmap patches are 2D grids where spatial patterns indicate geometric properties. Convolutional Neural Networks~\cite{lecun1998gradient} excel at extracting features from such grid-structured data through local receptive fields and weight sharing.

We employ a CNN to process heightmap patches. Early convolutional layers detect basic spatial patterns like edges and gradients. Deeper layers combine these into higher-level geometric structures such as corners, valleys, and plateaus. Finally, pooling aggregates the spatial information into a fixed-size embedding suitable for downstream processing.

% \begin{figure}[h]
% \centering
% \includegraphics[width=0.9\textwidth]{figures/heightmap_cnn.png}
% \caption{Heightmap CNN architecture processing patches to extract spatial features.}
% \label{fig:heightmap_cnn}
% \end{figure}

\subsection{Modeling Action Relationships: Set Transformer}

The number of valid actions varies dramatically---from dozens to hundreds depending on remaining items and available spaces. The value of placing item A in position 1 depends on whether item B has better alternatives. To address this interdependency, we need a mechanism for actions to reason about each other.

Transformer architectures~\cite{vaswani2017attention} allow each action to attend to all others through self-attention. We employ the Set Transformer~\cite{lee2019set}, specifically Induced Set Attention Blocks (ISAB), which achieve two critical properties:

\begin{enumerate}
    \item \textbf{Permutation Invariance}: The action set has no inherent order---action ordering is arbitrary. Set operations must respect this symmetry.
    \item \textbf{Computational Efficiency}: Standard attention has $\mathcal{O}(N^2)$ complexity where $N$ is the number of actions. ISABs use inducing points to reduce this to $\mathcal{O}(NM)$ where $M \ll N$, making attention tractable even when hundreds of actions are available.
\end{enumerate}

Through self-attention, each action's representation is contextualized by all other actions, enabling the network to reason about interdependencies, competition for resources, and relative action quality.

\subsection{Complete Architecture}

The complete Q-network integrates all components in a multi-stream architecture:

\begin{enumerate}
    \item \textbf{State Encoder}: Processes global state features (bin utilization, remaining items, constraints) through fully-connected layers to produce a state embedding
    \item \textbf{Action Encoder}: Processes action feature vectors through fully-connected layers to produce action embeddings
    \item \textbf{Heightmap Processing}: For each action, processes the corresponding heightmap patch through the CNN to extract spatial embeddings
    \item \textbf{Fusion}: Combines action, heightmap, and state embeddings for each candidate action
    \item \textbf{Attention}: Set Transformer processes all action embeddings jointly, producing contextualized representations
    \item \textbf{Q-Head}: Final layers map contextualized embeddings to Q-values $Q(s, a)$
\end{enumerate}

% \begin{figure}[h]
% \centering
% \includegraphics[width=\textwidth]{figures/network_architecture.png}
% \caption{Complete Q-network architecture showing information flow from state/action inputs through CNN, Set Transformer, to Q-values.}
% \label{fig:network_architecture}
% \end{figure}

This architecture extracts spatial patterns (CNN), reasons about action relationships (Transformer), and integrates global context to produce Q-values accounting for both local geometry and global optimization.

\subsection{Training Strategy}

\paragraph{Experience Replay} We use a replay buffer storing transitions $(s_t, a_t, r_t, s_{t+1}, \text{done}_t, \text{action\_mask}_{t+1})$. The action mask indicates valid actions in the next state, crucial for variable action spaces. Replay provides sample efficiency (learning from each transition multiple times) and decorrelates training data (breaking temporal correlations that can destabilize learning).

\paragraph{N-Step Returns} Standard Q-learning uses 1-step returns, which propagate rewards slowly. We encounter a credit assignment problem: when packing fails, which of the hundreds of placement decisions was responsible? To address this, we employ multi-step returns~\cite{sutton2018reinforcement}:

\begin{equation}
Q(s, a) \leftarrow \sum_{i=0}^{n-1} \gamma^i r_{t+i} + \gamma^n \max_{a'} Q(s_{t+n}, a')
\end{equation}

This propagates terminal rewards more rapidly through the episode. With 1-step returns, terminal reward requires hundreds of episodes to propagate back to early decisions. Multi-step returns accelerate credit assignment by incorporating multiple future rewards directly.

\paragraph{Target Networks} We use soft updates to the target network:
\begin{equation}
\theta_{\text{target}} \leftarrow \tau \theta_{\text{online}} + (1 - \tau) \theta_{\text{target}}
\end{equation}

This provides stable learning targets while allowing the target network to track the online network's improvements.

\paragraph{Exploration} We employ $\varepsilon$-greedy exploration with gradual decay, balancing exploration of new strategies with exploitation of learned policies.

\paragraph{Curriculum Learning} Training begins with simpler instances (fewer items, fewer constraints) and gradually increases difficulty~\cite{bengio2009curriculum}. This allows the agent to build foundational packing skills before tackling complex constraint combinations.

\section{Evolving the Architecture: Neurogenesis}

The neural architecture contains numerous hyperparameters: network depths, layer widths, attention heads, CNN channels, activation functions. We encounter a design challenge: manual tuning is time-consuming, requires domain expertise, and may miss better configurations in the vast hyperparameter space.

To address this, we turn to Neural Architecture Search (NAS)~\cite{elsken2019neural}, which automates the architecture design process. We employ an evolutionary approach inspired by neurogenesis---the growth of new neurons in biological networks~\cite{kempermann2015activity}.

\subsection{Genome Encoding}

Each candidate architecture is encoded as a genome---a structured representation of the hyperparameter configuration. The genome specifies architectural choices across all network components:

\begin{itemize}
    \item \textbf{Network Structure}: Encoder depths and widths, CNN channel configurations, attention mechanisms
    \item \textbf{Training Hyperparameters}: Learning rate, discount factor, temporal credit assignment horizon
    \item \textbf{Activation Functions}: Nonlinearity choices affecting gradient flow and expressiveness
\end{itemize}

The search space encompasses architectures ranging from small, fast networks to large, expressive ones, allowing the evolutionary process to discover the right balance for our problem.

\subsection{Evolutionary Process}

The evolutionary algorithm operates as follows:

\begin{enumerate}
    \item \textbf{Initialization}: Generate a population of $N$ random genomes to explore diverse regions of the architecture space
    \item \textbf{Evaluation}: Train each architecture for $K$ episodes, evaluate on validation set to measure generalization
    \item \textbf{Selection}: Select top $M$ performers as parents, preserving successful traits
    \item \textbf{Crossover}: Combine parent genomes to create offspring, inheriting complementary strengths
    \item \textbf{Mutation}: Randomly modify offspring with probability $p_{\text{mut}}$, introducing variation and exploring nearby configurations
    \item \textbf{Iteration}: Repeat for $G$ generations, progressively refining the population
\end{enumerate}

This process explores the architecture space efficiently. Crossover leverages successful configurations by combining proven components, while mutation discovers novel designs by perturbing existing ones~\cite{real2019regularized}. The evolutionary approach naturally balances exploration (trying new architectures) and exploitation (refining successful ones).

\subsection{Fitness Function}

Genomes are evaluated on:
\begin{itemize}
    \item \textbf{Packing performance}: Average bins used on validation instances
    \item \textbf{Training stability}: Variance in episode returns
    \item \textbf{Computational cost}: Parameter count and inference time
\end{itemize}

The fitness function balances these objectives, favoring architectures that pack efficiently without excessive complexity.

\section{The Complete System}

All components integrate into a complete learning system:

\begin{enumerate}
    \item Problem instances are loaded with items, bins, and constraints
    \item The EMS-based environment generates valid actions at each step
    \item Heightmap patches are extracted around each placement position
    \item The Q-network scores all actions based on features, heightmaps, and attention
    \item The highest-scoring valid action is executed
    \item Reward shaping provides learning signals
    \item Experience replay stores transitions for training
    \item Double DQN with n-step returns updates the Q-network
    \item Neurogenesis searches for optimal architectures
\end{enumerate}

% \begin{figure}[h]
% \centering
% \includegraphics[width=\textwidth]{figures/system_overview.png}
% \caption{Complete system architecture showing the flow from problem instance through EMS generation, neural network scoring, action selection, and learning loop.}
% \label{fig:system_overview}
% \end{figure}

This creates a system addressing the key challenges identified in earlier chapters:

\begin{itemize}
    \item \textbf{Constraint Satisfaction}: EMS provides collision-free candidates; environment validates complex constraints
    \item \textbf{Long-Term Optimization}: Reinforcement learning with shaped rewards optimizes cumulative packing quality
    \item \textbf{Spatial Reasoning}: Heightmap CNN extracts local geometric patterns
    \item \textbf{Action Relationships}: Set Transformer models interdependencies between placements
    \item \textbf{Credit Assignment}: N-step returns propagate terminal rewards efficiently
    \item \textbf{Architecture Design}: Neurogenesis discovers effective network configurations
\end{itemize}

Each design choice builds on insights from Chapter~\ref{ch:problem_analysis} and state-of-the-art methods reviewed in Chapter~\ref{ch:state_of_art}. The result is a system that learns effective packing strategies without labeled data, handles complex real-world constraints, and optimizes for long-term bin minimization.
